\questao{91}{corB}{Eletricidade}{Força elétrica (lei de Coulomb)}{Só de conceito}
\begin{enunciado}
% gerado por ferramentas/md2tex_enunciado.py a partir de questoes/q091.md (não editar à mão)
O modelo de Rutherford para o átomo consiste em um núcleo bastante denso, pequeno e carregado positivamente, ao redor do qual orbitam partículas negativas. A interação atrativa entre o núcleo e os elétrons foi explicada posteriormente pela existência de uma força centrípeta de origem elétrica, cujo módulo, segundo a lei de Coulomb, é proporcional à constante eletrostática e ao quadrado da carga elementar, e inversamente proporcional ao quadrado da distância entre o núcleo e os elétrons. No entanto, o modelo planetário proposto por Rutherford não concorda com a teoria eletromagnética clássica, segundo a qual os elétrons descreveriam um movimento em espiral e se chocariam com o núcleo, o que não corresponde à realidade.\par
\figq[0.24\linewidth]{q091_fig1.png}
Qual gráfico apresenta a força de atração entre o núcleo e o elétron, em função do tempo para o átomo de hidrogênio, levando em consideração a teoria eletromagnética clássica?\par
\begin{center}
\altfig{A}{q091_altA.png}\hfill\altfig{B}{q091_altB.png}\hfill\altfig{C}{q091_altC.png}\hfill\altfig{D}{q091_altD.png}\hfill\altfig{E}{q091_altE.png}
\end{center}

\end{enunciado}

\quer{qual gráfico mostra o que acontece com a força elétrica entre o núcleo e o elétron, com o passar do tempo, enquanto
o elétron ``cai'' em espiral na direção do núcleo.}

\begin{ideia}
Cargas de sinais opostos se atraem. O núcleo é positivo e o elétron é negativo, então um puxa o outro.
A lei de Coulomb diz de que tamanho é esse puxão:
\[ F = \frac{k \cdot e^2}{d^2} \]
Aqui, $k$ é uma constante e $e$ é a carga do elétron: \textbf{esses dois valores não mudam}. A única coisa que muda é a
distância $d$ entre o elétron e o núcleo. Como $d$ está \textbf{embaixo} da fração (e ao quadrado, isto é, $d \times d$), vale a regra:
\textbf{quanto menor a distância, maior a força, e ela cresce muito depressa.}

\textbf{Pense em um ímã de geladeira.} Longe da porta, você quase não sente nada. Chegando perto, o puxão aparece. Bem
pertinho, ele ``gruda'' de uma vez. A força não cresce aos poucos: ela dispara quando a distância fica pequena.
(A força do ímã é de outro tipo, mas se comporta de jeito parecido.)
\end{ideia}

\passoapasso
\npasso{Descubra o que acontece com a distância.}
Na figura, o elétron anda em espiral e vai se aproximando do núcleo. Então, com o passar do tempo, a distância $d$
\textbf{diminui}.

\npasso{Descubra o que acontece com a força.}
Distância menor quer dizer força maior. Então, com o passar do tempo, a força \textbf{aumenta}.
Só com isso já eliminamos três gráficos: o A (força sempre igual), o B e o E (força diminuindo).
Sobram o C e o D.

\npasso{A força cresce em linha reta (C) ou em curva (D)?}
O gráfico é de força contra \textbf{tempo}. Vamos testar com números fáceis. O valor de $k \cdot e^2$ não muda; para a conta
ficar simples, faça de conta que ele vale 36. Imagine que, a cada segundo, o elétron chega 1 passo mais perto
(distância 3, depois 2, depois 1):
\begin{conta}
\renewcommand{\arraystretch}{1.3}
\begin{tabular}{cccc}
\textbf{Tempo} & \textbf{Distância $d$} & \textbf{$d^2$} & \textbf{Força $= 36 \div d^2$}\\ \hline
0 & 3 & $3 \times 3 = 9$ & $36 \div 9 = 4$\\
1 & 2 & $2 \times 2 = 4$ & $36 \div 4 = 9$\\
2 & 1 & $1 \times 1 = 1$ & $36 \div 1 = 36$
\end{tabular}
\end{conta}
Os passos da distância são iguais (sempre 1), mas os saltos da força não: de 4 para 9 ela subiu 5; de 9 para 36 subiu 27.
\textbf{O aumento é cada vez maior.}

Um gráfico em linha reta (C) mostraria um aumento sempre igual. O gráfico que começa subindo devagar e depois dispara para
cima é o \textbf{D}. (Na teoria clássica, o elétron ainda cai cada vez mais depressa, o que deixa a subida final mais
forte ainda.)

\resposta{D}{A distância diminui com o tempo, então a força elétrica aumenta, e aumenta cada vez mais depressa.}

\begin{outras}
\alt{A} Força constante só aconteceria se a distância não mudasse (órbita circular). Na espiral, a distância muda.\par
\altl{B} e \alt{E} A força diminuiria se o elétron estivesse se \emph{afastando} do núcleo. Ele está se aproximando.\par
\alt{C} A força aumenta, mas não em linha reta: ela depende de $1/d^2$, e por isso cresce cada vez mais rápido.
\end{outras}

\begin{pegadinha}[Onde muita gente erra]
O texto diz que a força é ``inversamente proporcional ao quadrado da distância''. Quem lê depressa pensa: ``inversamente\ldots{}
então a força diminui'' e marca B ou E. Cuidado: \emph{inversamente} quer dizer que, quando uma coisa diminui, a outra aumenta.
Aqui a \textbf{distância diminui}; logo, a \textbf{força aumenta}.
\end{pegadinha}
